\documentclass[12pt]{article}
\usepackage[a4paper,margin=1in]{geometry}
\usepackage[utf8]{inputenc}
\usepackage[T1]{fontenc}
\usepackage{amsmath,amssymb}
\usepackage{textcomp}
\usepackage{parskip}
\begin{document}
\begin{center}
{\Large \textbf{2010 STEP II}}
\end{center}

\subsection*{Page 1}
\begin{flushleft}
\bigskip\textbf{Section A: Pure Mathematics}\par
\bigskip\textbf{Question 1.} Let P be a given point on a given curve C. The osculating circle to C at P is defined to be\par
the circle that satisfies the following two conditions at P : it touches C; and the rate of change\par
of its gradient is equal to the rate of change of the gradient of C.\par
Find the centre and radius of the osculating circle to the curve y = 1 - x + tan x at the point\par
\[on the curve with x-coordinate 1 \pi.\]
\bigskip\textbf{Question 2.} Prove that\par
\[cos 3x = 4 cos3 x - 3 cos x .\]
Find and prove a similar result for sin 3x in terms of sin x.\par
\textbf{(i)} Let\par
α\par
\[I(α) = 7 sin x - 8 sin3 x dx .\]
\[\int\]
0 ¡ ¢\par
Show that\par
\[I(α) = - 8 c^3 + c + 5 ,\]
3 3\par
\[where c = cos α. Write down one value of c for which I(α) = 0.\]
\textbf{(ii)} Useless Eustace believes that\par
sinn+1 x\par
\[sinn x dx =\]
Z n + 1\par
\[for n = 1, 2, 3, . . . . Show that Eustace would obtain the correct value of I(β) , where\]
\[cos β = - 1 .\]
Find all values of α for which he would obtain the correct value of I(α).\par
\bigskip\textbf{Question 3.} The first four terms of a sequence are given by F = 0, F = 1, F = 1 and F = 2. The\par
0 1 2 3\par
general term is given by\par
\[F = aλn + bµn , (∗)\]
n\par
where a, b, λ and µ are independent of n, and a is positive.\par
\textbf{(i)} Show that λ2 + λµ + µ2 = 2, and find the values of λ, µ, a and b.\par
\textbf{(ii)} Use (∗) to evaluate F .\par
\infty\par
F\par
n\par
\textbf{(iii)} Evaluate .\par
2n+1\par
X\par
\[n=0\]
\[9470 Jun10\]
\end{flushleft}
\subsection*{Page 2}
\begin{flushleft}
\bigskip\textbf{Question 4.} (i) Let\par
a f(x)\par
\[I = dx .\]
Z f(x) + f(a - x)\par
Use a substitution to show that\par
a f(a - x)\par
\[I = dx\]
Z f(x) + f(a - x)\par
and hence evaluate I in terms of a.\par
Use this result to evaluate the integrals\par
\[\pi\]
1 ln(x + 1) 2 sin x\par
dx and dx .\par
\[Z ln(2 + x - x^2) Z sin(x + \pi )\]
0 0 4\par
\textbf{(ii)} Evaluate\par
2 sin x\par
dx .\par
Z 1 x sin x + sin 1\par
2 x\par
¡ ¢\par
\bigskip\textbf{Question 5.} The points A and B have position vectors i + j + k and 5i - j - k, respectively, relative to the\par
origin O. Find cos 2α, where 2α is the angle ∠AOB.\par
\textbf{(i)} The line L has equation r = λ(mi + nj + pk). Given that L is inclined equally to OA\par
1 1\par
and to OB, determine a relationship between m, n and p. Find also values of m, n and\par
p for which L is the angle bisector of ∠AOB.\par
\textbf{(ii)} The line L has equation r = µ(ui + vj + wk). Given that L is inclined at an angle α\par
2 2\par
\[to OA, where 2α = ∠AOB, determine a relationship between u, v and w.\]
\[Hence describe the surface with Cartesian equation x^2 + y^2 + z^2 = 2(yz + zx + xy).\]
6 Each edge of the tetrahedron ABCD has unit length. The face ABC is horizontal, and P is\par
the point in ABC that is vertically below D.\par
\textbf{(i)} Find the length of P D.\par
\textbf{(ii)} Show that the cosine of the angle between adjacent faces of the tetrahedron is 1/3.\par
\textbf{(iii)} Find the radius of the largest sphere that can fit inside the tetrahedron.\par
\[9470 Jun10 [Turn over\]
\end{flushleft}
\subsection*{Page 3}
\begin{flushleft}
\bigskip\textbf{Question 7.} (i) By considering the positions of its turning points, show that the curve with equation\par
\[y = x^3 - 3qx - q(1 + q) ,\]
\[where q > 0 and q 6= 1, crosses the x-axis once only.\]
\textbf{(ii)} Given that x satisfies the cubic equation\par
\[x^3 - 3qx - q(1 + q) = 0 ,\]
and that\par
\[x = u + q/u ,\]
obtain a quadratic equation satisfied by u3. Hence find the real root of the cubic equation\par
\[in the case q > 0, q 6= 1.\]
\textbf{(iii)} The quadratic equation\par
\[t^2 - pt + q = 0\]
has roots α and β. Show that\par
\[α3 + β3 = p^3 - 3qp .\]
It is given that one of these roots is the square of the other. By considering the expression\par
\[(α2 - β)(β2 - α), find a relationship between p and q. Given further that q > 0, q 6= 1\]
and p is real, determine the value of p in terms of q.\par
\bigskip\textbf{Question 8.} The curves C and C are defined by\par
1 2\par
\[y = e\]
-x\par
\[(x > 0) and y = e\]
-x\par
\[sin x (x > 0),\]
respectively. Sketch roughly C and C on the same diagram.\par
1 2\par
Let x denote the x-coordinate of the nth point of contact between the two curves, where\par
n\par
\[0 < x < x < \cdots , and let A denote the area of the region enclosed by the two curves\]
1 2 n\par
between x and x . Show that\par
n n+1\par
\[A = 1 (e^2\pi - 1)e -(4n+1)\pi/2\]
n 2\par
\infty\par
and hence find A .\par
n\par
X\par
\[n=1\]
\[9470 Jun10\]
\end{flushleft}
\subsection*{Page 4}
\begin{flushleft}
\bigskip\textbf{Section B: Mechanics}\par
\bigskip\textbf{Question 9.} Two points A and B lie on horizontal ground. A particle P is projected from A towards B\par
at an acute angle of elevation α and simultaneously a particle P is projected from B towards\par
A at an acute angle of elevation β. Given that the two particles collide in the air a horizontal\par
distance b from B, and that the collision occurs after P has attained its maximum height h,\par
show that\par
\[2h cot β < b < 4h cot β\]
and\par
\[2h cot α < a < 4h cot α ,\]
where a is the horizontal distance from A to the point of collision.\par
10 (i) In an experiment, a particle A of mass m is at rest on a smooth horizontal table.\par
\[A particle B of mass bm, where b > 1, is projected along the table directly towards A\]
with speed u. The collision is perfectly elastic.\par
Find an expression for the speed of A after the collision in terms of b and u, and show\par
that, irrespective of the relative masses of the particles, A cannot be made to move at\par
twice the initial speed of B.\par
\textbf{(ii)} In a second experiment, a particle B is projected along the table directly towards A\par
with speed u. This time, particles B , B , . . . , B are at rest in order on the line\par
2 3 n\par
\[between B and A. The mass of B (i = 1, 2, . . . , n) is\]
λn+1-im,\par
\[where λ > 1. All\]
1 i\par
collisions are perfectly elastic. Show that, by choosing n sufficiently large, there is no\par
upper limit on the speed at which A can be made to move.\par
\[In the case λ = 4, determine the least value of n for which A moves at more than 20u.\]
You may use the approximation log 2 \approx 0.30103.\par
\bigskip\textbf{Question 11.} A uniform rod AB of length 4L and weight W is inclined at an angle θ to the horizontal. Its\par
lower end A rests on a fixed support and the rod is held in equilibrium by a string attached\par
to the rod at a point C which is 3L from A. The reaction of the support on the rod acts in a\par
direction α to AC and the string is inclined at an angle β to CA. Show that\par
\[cot α = 3 tan θ + 2 cot β .\]
\textasciicircum{}\circ \textasciicircum{}\circ \textasciicircum{}\circ\par
\[Given that θ = 30 and β = 45 , show that α = 15 .\]
\[9470 Jun10 [Turn over\]
\end{flushleft}
\subsection*{Page 5}
\begin{flushleft}
\bigskip\textbf{Section C: Probability and Statistics}\par
\bigskip\textbf{Question 12.} The continuous random variable X has probability density function f(x), where\par
\[a for 0 6 x < k\]
\par
\[f(x) = b for k 6 x 6 1\]
\par
\par
0 otherwise,\par
\par
\par
\[where a > b > 0 and 0 < k < 1. Show that a > 1 and b < 1.\]
\textbf{(i)} Show that\par
1 - 2b + ab\par
\[E(X) = .\]
2(a - b)\par
\textbf{(ii)} Show that the median, M , of X is given by M = if a + b > 2ab and obtain an\par
2a\par
expression for the median if a + b 6 2ab.\par
\textbf{(iii)} Show that M < E(X) .\par
13 Rosalind wants to join the Stepney Chess Club. In order to be accepted, she must play a\par
challenge match consisting of several games against Pardeep (the Club champion) and Quentin\par
(the Club secretary), in which she must win at least one game against each of Pardeep and\par
Quentin. From past experience, she knows that the probability of her winning a single game\par
against Pardeep is p and the probability of her winning a single game against Quentin is q,\par
\[where 0 < p < q < 1.\]
\textbf{(i)} The challenge match consists of three games. Before the match begins, Rosalind must\par
choose either to play Pardeep twice and Quentin once or to play Quentin twice and\par
Pardeep once. Show that she should choose to play Pardeep twice.\par
\textbf{(ii)} In order to ease the entry requirements, it is decided instead that the challenge match\par
will consist of four games. Now, before the match begins, Rosalind must choose whether\par
to play Pardeep three times and Quentin once (strategy 1), or to play Pardeep twice and\par
Quentin twice (strategy 2) or to play Pardeep once and Quentin three times (strategy 3).\par
\[Show that, if q - p > 1 , Rosalind should choose strategy 1.\]
\[If q - p < 1 give examples of values of p and q to show that strategy 2 can be better or\]
worse than strategy 1.\par
\[9470 Jun10\]
\end{flushleft}
\end{document}
